% Generated draft: requires mathematical and source review.
\documentclass[11pt]{article}
\usepackage[margin=25mm]{geometry}
\usepackage{fontspec}
\setmainfont{Latin Modern Roman}
\usepackage{amsmath,amssymb,mathtools,unicode-math}
\setmathfont{Latin Modern Math}
\usepackage{enumitem,needspace,xurl,hyperref}
\setlength{\parindent}{0pt}
\setlength{\parskip}{5pt}
\setlength{\emergencystretch}{3em}
\newcommand{\N}{\mathbb{N}}
\newcommand{\Z}{\mathbb{Z}}
\newcommand{\Q}{\mathbb{Q}}
\newcommand{\R}{\mathbb{R}}
\newcommand{\C}{\mathbb{C}}
\newcommand{\F}{\mathbb{F}}
\newcommand{\A}{\mathbb{A}}
\newcommand{\PP}{\mathbb{P}}
\newcommand{\E}{\mathbb{E}}
\newcommand{\eps}{\varepsilon}
\newcommand{\dd}{\,\mathrm{d}}
\begin{document}
\section*{UnsolvedMath ID: 30002906}
\textbf{Source problem code:} catalogue entry
\subsection{Definitions and mathematical statement}
The supplied record concerns a family of groups denoted by $\Gamma_g$ and subgroups $\Delta_g<\Gamma_g$ of finite index, meaning
\[
[\Gamma_g:\Delta_g]<\infty.
\]
Its title places the question in the setting of mapping class groups, but the record does not specify the underlying surfaces or define $\Gamma_g$. It also does not define the object $M_{\Delta_g}$, the cohomology coefficients, or the construction of the displayed map.

With these unresolved conventions retained, the source asks: for sufficiently natural finite-index subgroups $\Delta_g<\Gamma_g$, is the map
\[
H^k(M_{\Delta_g})\longrightarrow H^k(\Delta_g)
\]
surjective for $g\gg 0$? The source states the equivalent question as whether $H^k(\Delta_g)$ is pure.

Here ``sufficiently natural'' is an unspecified restriction on the permitted subgroups. The phrase $g\gg 0$ requires an eventual genus bound, but the record does not specify its dependence on $k$ or on the subgroup family, or whether $k$ is fixed before taking this bound. No weight structure or convention defining ``pure'' is supplied. Consequently, both the precise meaning of purity and the asserted equivalence with surjectivity require additional source context. This is an incompletely specified source question, not a fully determined mathematical assertion.
\subsection{Short English statement}
The question asks whether, in sufficiently large genus and for a suitably natural class of finite-index mapping class subgroups, the displayed cohomology map is onto, or equivalently, as asserted by the source, whether the subgroup cohomology is pure. The subgroup class, spaces, coefficients, map, purity convention, and quantifier order need clarification.
\subsection{Sources}
\begin{itemize}
\item Local UnsolvedMath registry, permanent ID $30002906$, title ``Cohomological Purity for Finite Index Mapping Class Subgroups,'' source problem code \texttt{catalogue entry}, in the Oberwolfach Reports Open Problems collection. The supplied proposed year is 2015; no individual proposer is identified.
\item Supplied report citation: \emph{New Perspectives on the Interplay between Discrete Groups in Low-Dimensional Topology and Arithmetic Lattices} (2016), Oberwolfach Reports, pp.~1701--1746, DOI \texttt{10.4171/owr/2015/30}. Supplied source URL and literature-source link, labelled ``original Oberwolfach report'': \url{https://doi.org/10.4171/owr/2015/30}. The record describes Oberwolfach Reports as documenting mathematical workshops and their open-problem discussions. No page locating this particular question is supplied.
\item Supplied literature assessment, checked on \texttt{2026-08-21}: the extracted homological purity question has no later settling theorem attached to the supplied record and is left unresolved there. The listed source checked is the original Oberwolfach report linked above. This is supplied dated curation, not an independently verified literature assessment; time-sensitive status and cited claims require independent verification.
\item The registry labels the problem ``open'' and the statement \texttt{corrected\_verified}. Its accompanying verification note reports only safe notation, typography, and TeX-delimiter cleanup, with hypotheses, quantifiers, and subquestions preserved. These supplied labels do not establish mathematical correctness or independent source verification.
\end{itemize}
\end{document}
